\documentclass[
    NAME={Dr. Helga Ingimundardóttir},
    EMAIL={helgaingim@hi.is},
    FACULTY={Iðnaðarverkfræði},
    TITLE={Sýnileg samvinna með tólum atvinnulífsins},
    SUBTITLE={PR vinnulag og gagnsæir ferlar í verkefnadrifinni verkfræðikennslu},
    SEMINAR={Kennsluráðstefna KHA},
    DATE={28. mars 2026},
    WIDE=true,
    ICELANDIC=true
]{HI-LaTeX//hi-beamer}
\usetikzlibrary{fit,matrix,backgrounds}
%\setbeameroption{show only notes}
\begin{document}



% ----------------------------------------------------
\begin{frame}{Af hverju þurfum við nýtt vinnulag?}
\begin{itemize}
  \item Háskólar vilja efla tengsl milli náms og starfs.
  \item Atvinnulífið krefst:
  \begin{itemize}
    \item gagnsæis,
    \item rekjanleika,
    \item fagmennsku í samvinnu og afhendingu.
  \end{itemize}
  \item Í verkefnadrifnum námskeiðum eru ferlar þó oft ósýnilegir.
  \item Þá verður erfitt að sjá:
  \begin{itemize}
    \item hver gerði hvað,
    \item hvernig ákvarðanir voru teknar,
    \item hvort endurgjöf var raunverulega nýtt.
  \end{itemize}
\end{itemize}


\note{
\begin{itemize}
  \item \emph{Hook}: þetta snýst um bilið milli \emph{náms og starfs}.
  \item Við metum oft \emph{niðurstöðu}, ekki \emph{ferlið}.
  \item Samvinna er stór hluti náms --- en er oft \emph{ósýnileg}.
  \item \textbf{Lykilsetning:} „Við sjáum hvað nemendur skila --- en ekki hvernig þeir komast þangað.“
\end{itemize}
}
\end{frame}

% ----------------------------------------------------
\begin{frame}{Vandinn í hefðbundinni teymisvinnu}
\begin{itemize}
  \item Framlag einstaklings verður oft óskýrt.
  \item Framvinda verkefnis er sjaldan sýnileg.
  \item Endurgjöf týnist milli lota eða skila.
  \item Námsmat byggir þá of mikið á lokaskjali eða lokakóða.
  \item Afleiðingar:
  \begin{itemize}
    \item laumufarþegar,
    \item misvægi í teymum,
    \item takmörkuð innsýn kennara í raunverulegt vinnuferli.
  \end{itemize}
\end{itemize}

\note{
\begin{itemize}
  \item Þetta eru \emph{þekkt vandamál}.
  \item Erfitt að sjá \emph{hver gerir hvað}.
  \item Endurgjöf verður \emph{einnota}.

  \item Sérstakt vandamál í \emph{teymisverkefnum}:
  \begin{itemize}
    \item nemendur \emph{skipta með sér verkum},
    \item „deila og drottna“ nálgun.
  \end{itemize}

  \item Spurningin er: \emph{eru allir að tileinka sér ÖLL hæfniviðmið námskeiðsins}?
  

  \item Kennari sér oft bara: \emph{lokaafurðina}, en ekki \emph{hver lærði hvað}.
  

  \item Þetta er \emph{áskorun fyrir réttlátt námsmat}.

  \item \textbf{Lykilsetning:} „Í teymisverkefnum vitum við oft hvað var gert --- en ekki hver lærði hvað.“
\end{itemize}
}

\end{frame}


% ----------------------------------------------------
\begin{frame}{PR-vinnulag sem náms- og matsinnviðir}
\begin{itemize}
  \item GitHub er leiðandi \emph{útgáfustýringarkerfi} í atvinnulífi.
  \item Pull request (PR) er tillaga að breytingu sem fer í rýni.
  \item PR tengir saman:
  \begin{itemize}
    \item breytingar,
    \item rökstuðning,
    \item jafningjarýni,
    \item viðbrögð við endurgjöf.
  \end{itemize}
  \item Endurgjöf verður hluti af skráðu ferli.
  \item Ferlið verður sýnilegt --- ekki bara niðurstaðan.
\end{itemize}

\note{
\begin{itemize}
  \item Ekki gera ráð fyrir \emph{forkunnáttu}. Nemendur fá \emph{ítarefni} og stuðning í byrjun.
  \item Þetta er \emph{vinnutæki} --- tekur smá tíma að læra.

  \item GitHub \emph{útgáfustýring}:
  \begin{itemize}
    \item staðall í \emph{atvinnulífi} (líka á Íslandi).
  \end{itemize}

  \item PR:
  \begin{itemize}
    \item leggur inn breytingu,
    \item fær \emph{rýni},
    \item svarar og lagar.
  \end{itemize}

  \item Nemendur átta sig fljótt á: þetta er \emph{hæfni} sem þau geta sett á \emph{CV}.
  Hjálpar að \emph{hvetja} þau að tileinka sér vinnulagið.
  


  \item \textbf{Lykilsetning:} „Þetta er ekki bara námsverkfæri --- þetta er færni sem þau taka með sér út á vinnumarkað.“
\end{itemize}
}

\end{frame}

% ----------------------------------------------------
\begin{frame}{Þrjár reglur sem skipta mestu máli}
\begin{enumerate}
  \item Skýr PR-lýsing með afmörkuðu umfangi.
  \item Allar efnislegar athugasemdir fá sýnileg svör.
  \item \emph{Merge} aðeins þegar breytingar hafa verið rýndar og samþykktar.
\end{enumerate}

\vspace{1em}
\begin{itemize}
  \item Þannig festist ábyrgð og fagmennska í vinnulaginu sjálfu.
\end{itemize}

\note{
\begin{itemize}
  \item Þetta eru \emph{grunnreglurnar}.
  \item Engar \emph{þöglar lagfæringar}.
  \item Nemendur þurfa að \emph{svara} --- ekki bara laga.

  \item Algeng \emph{tregða í byrjun}:    „Við ræddum þetta allt á Messenger“
  

  \item Vandinn:
  \begin{itemize}
    \item það er \emph{ekki sýnilegt},
    \item ekki \emph{rekjanlegt},
    \item kennari getur \emph{ekki metið það}.
  \end{itemize}

  \item Þess vegna:
  \begin{itemize}
    \item öll efnisleg umræða þarf að fara í \emph{PR}.
  \end{itemize}

  \item Þetta tekur smá tíma að \emph{venjast} --- en verður fljótt \emph{eðlilegt vinnulag}.

  \item \textbf{Lykilsetning:} „Ef það er ekki í PR --- þá er það ekki hluti af námsmatinu.“
\end{itemize}
}
\end{frame}


% ----------------------------------------------------
\begin{frame}{Tvö námskeið --- sama hugsun, ólík útfærsla}
\textbf{Upplýsingaverkfræði -- 2. ár BS}
\begin{itemize}
  \item Styttri lotur.
  \item Áhersla á að læra vinnuferlið sjálft.
  \item Nemendur fá skýran upphafs- og endapunkt í hverri lotu.
\end{itemize}

\vspace{0.75em}
\textbf{Viðskiptagreind -- 3. ár BS / 1. ár MS}
\begin{itemize}
  \item Sama geymsla allt misserið.
  \item Endurgjöf flyst áfram milli lota.
  \item Meiri áhersla á fagmennsku, rökstuðning og samskipti við fyrirtæki.
\end{itemize}



\note{
\begin{itemize}
  \item Sama \emph{hugmynd}, ólík \emph{dýpt}.
  \item Fyrra námskeið: \emph{þjálfun}.
  \item Seinna: \emph{raunverulegt vinnulag}.
  \item \textbf{Lykilsetning:} „Þetta fer frá æfingu --- yfir í fagmennsku.“
  \item Markmiðið sé ekki bara að nemendur skili „réttu svari“, heldur að þeir geti sýnt faglegt vinnulag.
\end{itemize}
}

\end{frame}



% ----------------------------------------------------
\begin{frame}{Dýpri fagmennska}
\begin{itemize}
  \item Raunveruleg gögn og samstarf.
  \item Sama gagnasafn allt misserið.
  \item Endurtekið ferli → dýpri skilningur.
  \item Nemendur læra að:
  \begin{itemize}
    \item rökstyðja,
    \item svara,
    \item afhenda faglega.
  \end{itemize}
\end{itemize}

\note{
\begin{itemize}
  \item Þetta er \emph{industry tengingin}.
  \item Ekki bara rétt svar --- heldur \emph{rétt vinnulag}.
  \item Sérstaklega mikilvægt þegar unnið er með ytri samstarfsaðila þar sem traust og skýr samskipti skipta máli.
  \item \textbf{Lykilsetning:} „Fagmennska verður sýnileg í hegðun.“
\end{itemize}
}

\end{frame}

% ----------------------------------------------------
\begin{frame}{Flæði námskeiðsins}
\begin{center}
    \begin{tikzpicture}
    [scale=0.5, transform shape,
    node distance=80mm,
    inner sep=3pt,
    auto,
    block/.style={
        draw,
        thick,
        fill=green!40,
        text width=38mm,
        align=center,
        minimum height=10mm
    },
    line/.style={
        draw, thick, ->, shorten >=2pt
    },
    dashedline/.style={
        draw, thick, ->, dashed, shorten >=2pt, color=black!50
    },
    week/.style={anchor=south west, text width=28mm},
    red/.style={block, draw=red, fill=red!40},
    orange/.style={block, draw=orange, fill=orange!40},
    yellow/.style={block, draw=yellow, fill=yellow!40},
    green/.style={block, draw=green, fill=green!40},
    violet/.style={block, draw=violet, fill=violet!40},
    blue/.style={block, draw=blue!40, fill=blue!20},
    agile/.style={block, draw=gray, fill=gray!10,inner sep=3mm}
    ]

    % Samantektarverkefni
    \node[week] (themecapstone) at (0,-75mm) {Vika 11--14  Capstone-verkefni};
    \node[blue, right of=themecapstone, xshift=70mm, label=below:\textbf{Verkefni}] (capstone) {Teymi (tCAP)};
    % Endurgjöf milli teyma
    \draw [line] (capstone) edge[loop above] (capstone);
    \node[blue, right of=capstone, label=below:\textbf{Ígrundun}] (capref) {Nemandi (iCAP)};
    \draw[line] (capstone) to (capref);

    % Litur/lotur
    \foreach \week/\theme/\color/\pos in {1--2/Gagnasiðfræði/red/1, 3--4 /Mállíkön/orange/2, 5--6/Leiðbeint nám/yellow/3, 7--8 /Klösun/green/4, 9--10/Ferlagreining/violet/5}{
        \node[week] (week\pos) at (0,-\pos*10mm) {Vika \week\\\theme};
        \node[\color, right of=week\pos, xshift=-50+10*\pos,label=below:\textbf{Undirbúningur}] (RAT\pos)  {Nemandi (iRAT)\\Teymi (tRAT)};
        \node[\color, right of=RAT\pos,label=below:\textbf{Verkefni}] (APP\pos) {Teymi (tAPP)};
        \node[\color, right of=APP\pos,label=below:\textbf{Ígrundun}] (REF\pos) {Nemandi (iREF)};

        % Línur milli eininga
        \path [line] (RAT\pos) -- (APP\pos);
        \path [line] (APP\pos) -- (REF\pos);

        % Jafningjamat innan teymis
        \draw [line] (APP\pos) edge[loop above] (APP\pos);
        
        % Tengingar við samantektarverkefni
        \path [dashedline] (APP\pos.east) to[out=-20, in=70] (capstone.north east);
        \path [dashedline] (REF\pos.east) to[out=-30, in=40] (capstone.east);
    }
    \begin{scope}[on background layer]
    \node[agile,fit=(RAT1) (RAT5),label=above:{\textbf{Áætlun spretts (sprint planning)}}] () {};
    \node[agile,fit=(APP1) (APP5),label=above:{\textbf{Sprettur (sprint)}}] () {};
    \node[agile,fit=(REF1) (REF5),label=above:{\textbf{Verkefnalisti (backlog)}}] () {};
    \end{scope}
\end{tikzpicture}
\end{center}

\note{
\begin{itemize}
  \item Viðskiptagreind byggir á byggi á endurteknum lotum þar sem vinna, endurgjöf, rýni og umbætur tengjast saman.
  \item Ekki ein lokaskil --- heldur samfellt \emph{ferli}.
  \item \textbf{Lykilsetning:} „Námið gerist í ferlinu --- ekki bara í lokaskilunum.“
\end{itemize}
}
\end{frame}


% ----------------------------------------------------
\begin{frame}{Hvað breytist?}
\textbf{Fyrir nemendur}
\begin{itemize}
  \item Sýnilegra framlag í teymisvinnu.
  \item Skýrari væntingar og ábyrgð.
  \item Minni laumufarþegavandi.
  \item Vinnulag sem flyst beint yfir í atvinnulíf.
\end{itemize}

\vspace{0.75em}
\textbf{Fyrir kennara}
\begin{itemize}
  \item Betri innsýn í framvindu og einstaklingsframlög.
  \item Endurgjöf varðveitist og nýtist áfram.
  \item Skilvirkara námsmat og minna þarf að lesa á milli línanna.
\end{itemize}

\note{
\begin{itemize}
  \item Þetta er \emph{praktíski ávinningurinn}.
  \item Ávinningur fyrir bæði nemendur og kennara --- þetta er ekki bara „meira eftirlit“, heldur betra nám.
  \item \textbf{Lykilsetning:} „Þetta gerir ósýnilega vinnu sýnilega.“
\end{itemize}
}
\end{frame}


% ----------------------------------------------------
\begin{frame}{Lykilatriði}
\begin{itemize}
  \item Góð teymiskennsla þarf sýnilega ferla, ekki bara lokaafurðir.
  \item Fagmennska er hegðun sem þarf að þjálfa markvisst.
  \item Tól atvinnulífsins geta orðið hluti af námsmati --- ekki bara aukabúnaður.
\end{itemize}

\vspace{1em}
\begin{center}
\large
PR-vinnulag færir námsmat nær raunverulegu starfi.
\end{center}


\note{
\begin{itemize}
  \item Halda þessu \emph{einföldu}.
  \item \textbf{Lykilsetning:} „Námsmat þarf að endurspegla hvernig unnið er.“
\end{itemize}
}
\end{frame}

% ----------------------------------------------------
\begin{frame}{Takk fyrir}
\Large
{\centering
\textbf{Spurningar?}\\[18pt]
}

\small
\textbf{Birtingar:}
\begin{itemize}
  \item Ingimundardóttir, H (2024). \emph{Endurskoðun á námskeiði í Viðskiptagreind: Hagnýt hæfni í brennidepli}. Ráðstefna Kennsluakademíu opinberu háskólanna, Veröld Háskóla Íslands.

  \item Ingimundardóttir, H. (2025). \emph{Viðskiptagreind sem brú milli náms og starfs}. 
  Tímarit Kennslumiðstöðvar Háskóla Íslands, 12(1).

  \item Ingimundardóttir, H. (2026). \emph{Designing authentic industry-engaged assessment for professional competence in business intelligence}. 
  22nd CDIO International Conference, University of Liverpool (samþykkt).
  
  \item Ingimundardóttir, H. (2026). \emph{Using pull requests to make collaboration visible in CDIO project-based courses}. 
  22nd CDIO International Conference, University of Liverpool (samþykkt).


\end{itemize}

\end{frame}

\end{document}

\end{document}